% !PRINTED 497
\weberpage{497}

Das Product dieser homogenen Functionen kann nach dem oben Bewiesenen nicht verschwinden, wenn keiner der Factoren verschwindet, und sein Grad ist gleich der Summe der Grade der einzelnen Factoren.

Die Zahlen des K\"orpers $\Omega$ sind unter den Functionen mit enthalten. Man erh\"alt sie, wenn man entweder den Grad oder die Anzahl der Variablen auf Null heruntersinken l\"asst.

Ueber die ganzen Functionen in $\Omega$ haben wir im ersten Bande mehrere wichtige S\"atze abgeleitet, von denen wir hier den Satz \S.\,141, IV hervorheben:

1. \gesperrt{Unter den ganzen Functionen in $\Omega$ m\"ussen reducible und irreducible unterschieden werden, von denen die ersten als Product aus mehreren ganzen Functionen in $\Omega$ darstellbar sind, die anderen nicht. Die reduciblen Functionen $\varphi$ lassen sich in eine endliche Anzahl von irreduciblen Factoren zerlegen, die selbst ganze Functionen in $\Omega$ sind, und die irreduciblen Factoren von $\varphi$ sind, von constanten Factoren, d. h. Zahlen in $\Omega$, abgesehen, eindeutig durch $\varphi$ selbst bestimmt.}

Als specielle F\"alle sind unter den ganzen Functionen in $\Omega$ auch die ganzen Functionen im K\"orper der rationalen Zahlen $R$ enthalten:
\weq{2}{\Phi(x,y,z,\ldots) = \sum a\,x^r y^s z^t \ldots,}
worin die Co\"efficienten $a$ rationale Zahlen sind.

Wenn diese Co\"efficienten \gesperrt{ganze Zahlen ohne gemeinsamen Theiler} sind, so heisst diese Function eine \gesperrt{urspr\"ungliche oder primitive}, von denen wir im Bd. I, \S.\,2 den Satz nachgewiesen haben:

2. \gesperrt{Das Product von zwei primitiven Functionen ist wieder eine primitive Function.}

Wenn die Co\"efficienten der Function (2), die wir f\"ur den Augenblick mit
\[ a_0, a_1, a_2, \ldots \]
bezeichnen wollen, ganze Zahlen mit dem gr\"ossten gemeinschaftlichen Theiler $m$ sind, so ist, wenn $m>1$ ist, $\Phi$ eine imprimitive \gesperrt{ganze ganzzahlige Function vom Theiler $m$, % [sic] \textquotedbl{}ganze ganzzahlige\textquotedbl{} so gedruckt
und der Theiler einer primitiven Function ist $= 1$.}
